% ============================
% LaTeX Template - Problem Set
% ============================
\documentclass[11pt]{article}

% ---------- Page + layout ----------
\usepackage[margin=1in]{geometry}
\usepackage{setspace}
\setstretch{1.1}
\usepackage{parskip}

% ---------- Math ----------
\usepackage{amsmath, amssymb, amsfonts, amsthm, mathtools}

% ---------- Tables ----------
\usepackage{booktabs}

% ---------- Links ----------
\usepackage[hidelinks]{hyperref}

% ---------- Header / footer ----------
\usepackage{fancyhdr}
\pagestyle{fancy}
\fancyhf{}
\lhead{\textbf{\course}}
\chead{\textbf{\pset}}
\rhead{\textbf{\name}}
\cfoot{\thepage}

% =========================
% Metadata (edit if needed)
% =========================
\newcommand{\name}{Tate Mason}
\newcommand{\course}{Econometrics II}
\newcommand{\pset}{Problem Set 4}
\newcommand{\duedate}{March 6, 2026}

% =========================
% Question / solution macros
% =========================
\newcounter{q}
\renewcommand{\theq}{\arabic{q}}

\newcommand{\question}[1]{%
  \refstepcounter{q}
  \vspace{0.75em}
  \noindent{\Large\textbf{Question \theq.} \normalsize\textbf{#1}}%
  \vspace{0.25em}
  \par
}

\newcounter{subq}[q]
\renewcommand{\thesubq}{\alph{subq}}

\newcommand{\subquestion}[1]{%
  \refstepcounter{subq}
  \vspace{0.35em}
  \noindent\textbf{(\thesubq)} \textbf{#1}\par
}

\newenvironment{solution}{%
  \vspace{0.15em}
  \noindent\textit{Solution. }\ignorespaces
}{\vspace{0.5em}\par}

% =========================
% Easy math macros
% =========================
\newcommand{\R}{\mathbb{R}}
\newcommand{\N}{\mathbb{N}}
\DeclareMathOperator{\E}{\mathbb{E}}
\DeclareMathOperator{\Var}{Var}
\DeclareMathOperator{\Cov}{Cov}
\DeclareMathOperator{\plim}{plim}
\DeclareMathOperator*{\argmin}{argmin}
\DeclareMathOperator*{\argmax}{argmax}
\newcommand{\P}{\mathbb{P}}

\newcommand{\vct}[1]{\mathbf{#1}}
\newcommand{\mtx}[1]{\mathbf{#1}}
\newcommand{\dd}{\,\mathrm{d}}
\newcommand{\eps}{\varepsilon}
\newcommand{\1}{\mathbf{1}}

% =========================
% Document
% =========================
\begin{document}

% ---------- Title block ----------
{\Large\textbf{\course} \hfill \textbf{\pset}}\\
\textbf{Name:} \name \hfill \textbf{Due:} \duedate\\
\rule{\textwidth}{0.4pt}
\question{Use “dataHW4 Problem1.mat” to complete this exercise. In this problem you will work with a
multinomial choice model very similar to the one from previous problem sets. Upon graduating
high school, individuals have three options: work (Y=0), attend a 2-year college (Y=1), attend a
4-year college (Y=2). The available student demographic characteristics (Xi) are a constant, an
indicator for whether either parent graduated from a 4-year college, and the student’s high school
GPA respectively. There is one choice specific variable, price measured in thousands of dollars
(Zprice). Each column of Zprice reflects the price for the corresponding choice. In other words, the
first column is the price for the work choice, the second column for the 2-year choice, and so on.
Note that I assume that price for the work option is zero.}

\subquestion{Estimate a mixed logit using simulation with R = 500 draws. The coefficient on price is nor-
mally distributed with mean µγ and variance σ2γ. The unknown coefficients for the student
demographic variables do not vary in the population. In other words, the only random coeffi-
cient is on price. Normalize the utility of work to zero. Report parameter estimates and SEs
for the βs, µγ, and σγ.}

\begin{solution}
  \begin{table}[htbp]
  \centering
    \caption{Mixed Logit Estimates}
    \begin{tabular}{lcc}
      \hline
      & Parameter & S.E. \\
      \hline
      \multicolumn{3}{l}{\textit{2-Yr}} \\
      $\beta_{0,1}$    & -0.930576 & (0.377443) \\
      $\beta_{BA,1}$   &  0.210334 &  (0.113812) \\
      $\beta_{GPA,1}$  &  0.416573 &  (0.122459) \\

      \multicolumn{3}{l}{\textit{4-Yr}} \\
      $\beta_{0,2}$    & -2.315272 &  (1.022829) \\
      $\beta_{BA,2}$   & 0.701975 &  (0.239203) \\
      $\beta_{GPA,2}$  & 1.176881 &  (0.295022) \\
      \hline
      $\mu_\gamma$     & -0.148575 &  (0.008202) \\
      $\sigma_\gamma$  & 0.015835 &  (1.587908) \\
      \hline
    \end{tabular}
  \end{table}
\end{solution}

\subquestion{Calculate the average marginal effect of having a parent with a BA on the likelihood of attending
a 4-year college. To calculate the average marginal effect you will need to integrate out over
the price coefficient, similar to what you did in estimation. Here you can use 100 draws per
observation. Then bootstrap standard error.}

\begin{solution}
  \centering
  \caption{Average Marginal Effect of Parent Education on 4-Year College Attendance}

  \begin{tabular}{lcc}
  \hline
    AME & Standard Error \\
  \hline
    0.111298 & 0.028921 \\
  \hline
  \end{tabular}

\end{solution}

\subquestion{Re-estimate the mixed logit model from part (a) using 5-point Gaussian-Hermite quadrature
to approximate the integral over the choice probabilities.}

\begin{solution}
  \begin{table}[htbp]
    \centering
    \caption{Mixed Logit Estimates (Gaussian-Hermite Quadrature)}
    \begin{tabular}{lcc}
      \hline
      & Parameter & S.E. \\
      \hline
      \multicolumn{3}{l}{\textit{2-Yr}}
      $\beta_{0,1}$    & -0.919686 & (0.072787) \\
      $\beta_{BA,1}$   &  0.217438 &  (0.039067) \\
      $\beta_{GPA,1}$  &  0.428476 &  (0.021781) \\
      \multicolumn{3}{l}{\textit{4-Yr}} \\
      $\beta_{0,2}$    & -2.348588 &  (0.111559) \\
      $\beta_{BA,2}$   & 0.719959 &  (0.040966) \\
      $\beta_{GPA,2}$  & 1.208237 &  (0.02856) \\
      \hline
      $\mu_\gamma$     & -0.155445 &  (0.006060) \\
      $\sigma_\gamma$  & 0.038863 &  (0.093009) \\
      \hline
    \end{tabular}
  \end{table}
\end{solution}

\question{Use “dataHW4 Problem2.mat” to complete this exercise. In this problem you will work with panel
data to estimate a logit model that allows for discrete unobserved heterogeneity. During the pan-
demic working from home became the norm, and many people continue to work from home. Over
the course of six months, you collect data from women between the ages of 30 and 40 who are work-
ing full-time. Each month you observe whether the woman was working at home full time (Y=1).
The six columns of Y reflect her response for each month. For each respondent you also observe Xi
containing the number of children at home and years of schooling (in addition to a constant). You
also observe distance from the physical workplace in miles (Zdist). Note that Xi and distance are
fixed for each women over the 6 months.}

\subquestion{You suspect that in addition to children and schooling, women likely differ in their tastes for
working at home. Assume that there are two unobserved types of women, one type prefers the
office, the other prefers to work from home, all else equal. Type is assumed to be uncorrelated
with Xi or distance. Estimate a mixed logit model for the decision to work from home that allows
for these two types of women. Note that this means there will be two additional parameters
to estimate beyond the coefficients associated with Xi and distance. You need to estimate the
proportion of Type 2 women, and the additional preference these women have for working at
home. Report coefficients and SEs.}

\begin{solution}
  \begin{table}[htbp]
    \centering
    \caption{Mixed Logit Estimates (Two Types of Women)/EM Algorithm}
    \begin{tabular}{lcc}
      \hline
      & Parameter & S.E. \\
      \hline
      \multicolumn{3}{l}{\textit{Likelihood Estimates}} \\
      $\beta_{0}$    & -0.468813 & (0.181709) \\
      $\beta_{Children}$   &  0.511097 &  (0.011169) \\
      $\beta_{Schooling}$  &  -0.204464 &  (0.007455) \\
      $\gamma$ & 0.052225 & (0.000165) \\
      $\delta$ & 1.98194 & (0.068723) \\
      $\pi_2$ & 0.680251 & (0.020931) \\
      \hline
      \multicolumn{3}{l}{\textit{EM Algorithm Estimates}} \\
      $\beta_{0}$    & -0.468813 & --- \\
      $\beta_{Children}$   &  0.511097 &  --- \\
      $\beta_{Schooling}$  &  -0.204464 &  --- \
      $\gamma$ & 0.052225 & --- \\
      $\delta$ & 1.98194 & --- \\
      $\pi_2$ & 0.680251 & --- \\
      \hline
    \end{tabular}
  \end{table}

The standard errors from the EM algorithm could be obtained through the numerical Hessian, taking the inverse of the Hessian at the estimated values provides the covariance matrix. From there we can take the root of the diagonal elements.
If this is too computationally taxing, bootstrapped standard errors can be obtained by running the EM on $B$ iterations of the data and calculating the standard deviation of the parameter estimates across the $B$ iterations.
\end{solution}



\end{document}
